\documentclass[12pt,a4paper]{article}
\usepackage[a4paper,margin=20mm]{geometry}
\usepackage{amsmath}
\usepackage{amssymb}

\begin{document}

\begin{center}
  {\Large\bfseries Answer Key}
\end{center}

\begin{enumerate}
  \item \textbf{Sometimes true.} True for \(x=3\); false for \(x=0.5\) (counterexample).

  \item \textbf{Sometimes true.} True for positive \(x\); false for \(x=0\) and for negative \(x\) (e.g. \(x=-2\)).

  \item
  (a) \textbf{False}: \(4x+y\) has no constant term.\\
  (b) \textbf{False}: \(2a+3b+1\) has \(3\) terms but only \(2\) variables.\\
  (c) \textbf{False}: if \(x=-2\), then \(-3x=6\), which is not negative.\\
  (d) \textbf{False}: \(x\) has coefficient \(1\), which is not bigger than \(1\).

  \item
  (a) \textbf{Always true}.\\
  (b) \textbf{Sometimes true}.\\
  (c) \textbf{Never true}.\\
  (d) \textbf{Sometimes true}.

  \item Sam is wrong: the statement is \textbf{sometimes true}. Counterexample: \(x=0\), since
  \[
    3x=0,\qquad x+1=1,
  \]
  so \(3x\) is not greater than \(x+1\).

  \item Example: \(x+2>5\) is \textbf{sometimes true}. It is true for \(x=4\) and false for \(x=1\).
\end{enumerate}

\end{document}